\subsection{A sixth-root specialization in every position}
\label{one2:sec:sixthroot}

The special role of sixth roots in relations for Nielsen polylogarithms
is studied by Borwein and Straub~\cite{BorweinStraub2015}.  The preceding
position formula gives a direct specialization in that established
setting, including an explicit rational coefficient for one component
at every weight.  Throughout this subsection,
\[
 \omega=e^{i\pi/3},\qquad N=a+b+2,\qquad a,b\ge0,
\]
and all logarithms and polylogarithms use their principal branches.

\begin{corollary}[Sixth-root closure and a pure parity component]
\label{one2:cor:sixthroot}
At $z=\omega$, formula~\eqref{one2:eq:closed} reduces
$U_{a,b}(\omega)$ to classical values $\Li_k(\omega)$, zeta values,
and powers of $i\pi$, since
\begin{equation}
 \frac1{1-\omega}=\omega,\qquad
 -\Log(1-\omega)=\Log\omega=\frac{i\pi}{3}.
 \label{one2:eq:sixthfixed}
\end{equation}
Moreover,
\begin{equation}
 U_{a,b}(\omega)+(-1)^N\overline{U_{a,b}(\omega)}
       =R_{N,a}(i\pi)^N,\qquad R_{N,a}\in\mathbb Q.
 \label{one2:eq:sixthparity}
\end{equation}
Thus $\Re U_{a,b}(\omega)$ is a rational multiple of $\pi^N$
when $N$ is even, and $\Im U_{a,b}(\omega)$ is a rational multiple
of $\pi^N$ when $N$ is odd.  Explicitly, with Bernoulli polynomials
defined by
\[
 \frac{t e^{xt}}{e^t-1}
    =\sum_{k\ge0}B_k(x)\frac{t^k}{k!},\qquad B_k=B_k(0),
\]
one may take
\begin{align}
 R_{N,a}
 &=-\frac{2}{3^N N!}
 -\sum_{k=a+1}^{N}
   \frac{(-1)^{a+k}\binom{k-1}{a}\,2^k B_k(1/6)}
        {3^{N-k}(N-k)!k!}\notag\\
 &\quad+(-1)^{N-a}
   \sum_{\substack{N-a\le m\le N\\m\ \mathrm{even}}}
   \frac{\binom{m-1}{N-a-1}\,2^m B_m}
        {3^{N-m}(N-m)!m!}.
 \label{one2:eq:sixthcoefficient}
\end{align}
The same formulas at $\overline\omega$ follow by complex conjugation.
\end{corollary}
\begin{proof}
The identities in \eqref{one2:eq:sixthfixed} follow from
$1-\omega=\overline\omega$.  Along the radial segment
$z=r\omega$, $0\le r\le1$, the real part of $1-z$ is at least
$1/2$, so the indicated logarithm has no branch crossing.
The transformed path lies in the upper half-plane for $r>0$,
starts at $1$, and ends at $\omega$; the endpoint limits already
used to prove \eqref{one2:eq:closed} apply.  Consequently that
formula has exactly the stated principal values.  The ordinary
boundary multiple series agree with their radial limits by the
same summation-by-parts argument used at $z=i$.

Put $T=i\pi$ and $L=T/3$.  The classical Bernoulli Fourier identity is
\begin{equation}
 \Li_k(e^{2\pi i x})+(-1)^k\Li_k(e^{-2\pi i x})
       =-\frac{(2\pi i)^k}{k!}B_k(x),
 \qquad 0<x<1,\quad k\ge1.
 \label{one2:eq:bernoullifourier}
\end{equation}
For $k=1$, this follows directly from $\Li_1(z)=-\Log(1-z)$
on the stated arc.  Successive integration, with the zero average
of both sides fixing the constants, proves every higher case.
Since $\overline L=-L$, the parity projection of a classical term
in \eqref{one2:eq:closed} is
\begin{align*}
 &L^{N-k}\Li_k(\omega)
       +(-1)^N\overline{L^{N-k}\Li_k(\omega)}\\
 &\hspace{1cm}=L^{N-k}
    \bigl(\Li_k(\omega)+(-1)^k\Li_k(\overline\omega)\bigr)
   =-\frac{2^k B_k(1/6)}{3^{N-k}k!}\,T^N.
\end{align*}
The polynomial term $-L^N/N!$ contributes
$-2T^N/(3^N N!)$.  A zeta term with zeta weight $m$ contributes
zero when $m$ is odd.  For even $m\ge2$, use
\[
 \frac{2\zeta(m)}{(i\pi)^m}=-\frac{2^m B_m}{m!}.
\]
In the last sum of \eqref{one2:eq:closed}, write
$m=b+j+2=N-a+j$.  Its binomial coefficient becomes
$\binom{m-1}{N-a-1}$, while its sign after the even-zeta evaluation
is $(-1)^b=(-1)^{N-a}$.  Combining these three contributions proves
\eqref{one2:eq:sixthcoefficient} and hence
\eqref{one2:eq:sixthparity}.
\end{proof}

There is a precise geometric reason that this specialization closes
on a unit-circle argument.  If $|z|=1$ and $z\ne1$, then
\[
 \left|\frac1{1-z}\right|=1
 \quad\Longleftrightarrow\quad
 |1-z|^2=2-2\Re z=1
 \quad\Longleftrightarrow\quad z\in\{\omega,\overline\omega\}.
\]
Both of these points are fixed by $z\mapsto(1-z)^{-1}$.
This classifies the unit-circle endpoint geometry of this particular
transformation.

\begin{corollary}[The two weight-three values]
\label{one2:ex:sixththree}
Define the ordinary Clausen value by its absolutely convergent series
\[
 C=\operatorname{Cl}_2(\pi/3)
   =\sum_{n\ge1}\frac{\sin(n\pi/3)}{n^2}.
\]
Then
\begin{align}
 U_{0,1}(\omega)=\Li_{2,1}(\omega)
  &=\frac23\zeta(3)-\frac\pi3 C+\frac{i\pi^3}{324},
    \label{one2:eq:sixth21}\\
 U_{1,0}(\omega)=\Li_{1,2}(\omega)
  &=-\frac43\zeta(3)+\frac\pi3 C+\frac{i\pi^3}{324}.
    \label{one2:eq:sixth12}
\end{align}
\end{corollary}
\begin{proof}
The classical values needed are
\[
 \Li_1(\omega)=\frac{i\pi}{3},\qquad
 \Li_2(\omega)=\frac{\pi^2}{36}+iC,\qquad
 \Li_3(\omega)=\frac{\zeta(3)}3+\frac{5i\pi^3}{162}.
\]
The real dilogarithm component and imaginary trilogarithm component
follow from \eqref{one2:eq:bernoullifourier}.  The real
trilogarithm component follows by grouping the defining series
into residue classes, or from the distribution identity
\[
 \Li_s(\omega)+\Li_s(\overline\omega)
   =(1-2^{1-s})(1-3^{1-s})\zeta(s),\qquad s>1.
\]
For example, summing over the three cube roots gives
$(3^{1-s}-1)\zeta(s)$ for the two nontrivial cube roots;
applying the two-term distribution formula to the pair of sixth
roots multiplies that sum by $2^{1-s}-1$.
Now \eqref{one2:eq:height} at $p=2$ gives, with $L=i\pi/3$,
\[
 U_{0,1}(\omega)
   =\zeta(3)-\Li_3(\omega)+L\Li_2(\omega)-\frac23L^3.
\]
Substitution yields \eqref{one2:eq:sixth21}.  The shuffle identity
\[
 \Li_1(\omega)\Li_2(\omega)
     =U_{1,0}(\omega)+2U_{0,1}(\omega)
\]
then proves \eqref{one2:eq:sixth12}.  In both positions,
\eqref{one2:eq:sixthcoefficient} gives $R_{3,a}=-1/162$.
\end{proof}

The independent script \texttt{code/independent/one\_two/sixth\_root.py}
computes \eqref{one2:eq:sixthcoefficient} with exact rational
arithmetic and compares the specialized closed formula, its pure
parity component, and the displayed weight-three examples with the
logarithmic moment.  It covers every position of the $2$ at weights
$2$ through $8$, at both $\omega$ and $\overline\omega$, using
$100$ working decimal digits.  Across these $56$ cases, the largest
closed-form versus moment residual is below $6.59\cdot10^{-100}$;
the largest parity-projection residual is below $7.15\cdot10^{-102}$.
The four checks of the two explicit examples and their conjugates
have residual below $2.16\cdot10^{-101}$.  These are numerical checks
of the proved identities, rather than rigorous quadrature error bounds.
